\section{Starten des Programms}
\subsection{Abhängigkeiten installieren}
Im Ordner \texttt{SkylineApp} ein Terminal bzw.\ eine Eingabeaufforderung öffnen und ausführen:
\begin{lstlisting}[language=sh,basicstyle=\ttfamily\small]
npm install
\end{lstlisting}
Dies installiert alle benötigten Pakete (kann einige Minuten dauern).

\subsection{Entwicklungsserver starten}
Anschließend den Expo-Entwicklungsserver starten:
\begin{lstlisting}[language=sh,basicstyle=\ttfamily\small]
npm start
\end{lstlisting}
Es erscheint ein QR-Code und die Metro-Bundler-Oberfläche.

\subsection{App auf dem Gerät öffnen}
\begin{itemize}
  \item \textbf{Smartphone:} Mit der Kamera oder der \emph{Expo Go}-App den angezeigten QR-Code scannen. Rechner und Smartphone müssen im gleichen WLAN sein.
  \item \textbf{Android-Emulator:} \texttt{a} im Terminal drücken.
  \item \textbf{iOS-Simulator (nur macOS):} \texttt{i} im Terminal drücken.
\end{itemize}

Bei Netzwerkproblemen (z.\,B.\ wenn sich Rechner und Smartphone nicht im selben Netz befinden) kann der Tunnel-Modus verwendet werden: \texttt{npx expo start --tunnel}

\subsection{Alternative: App aus dem iOS App Store}
Skyline ist im Apple App Store veröffentlicht. Die App kann dort kostenlos heruntergeladen und getestet werden, ohne die obigen Schritte auf einem Rechner durchzuführen.
